\documentclass[11pt,a4paper]{article}
\usepackage[margin=1in]{geometry}
\usepackage{longtable}
\usepackage{booktabs}
\usepackage{array}
\usepackage{hyperref}
\usepackage{pdflscape} % For landscape pages
\usepackage{todonotes}
\usepackage{comment}

\makeatletter
\renewcommand\paragraph{%
  \@startsection{paragraph}{4}{\z@}% <-- 3: indent (set to 0)
    {4mm}% <-- 4: Space _before_
    {-1em}% <-- 5: Space _after_ (negative for run-in title)
    {\normalfont\normalsize\bfseries$\bullet$~}% 6: Style
}
\makeatother


% Define a new column type for the tables to allow wrapping text
% CORRECTED: Added [1] to indicate it takes one argument (the width)
\newcolumntype{L}[1]{>{\raggedright\arraybackslash}p{#1}}

\title{Strategic Imperative: Fully Integrating the Kumar Lab to Accelerate Envision's 2026 Roadmap}
\author{}
\date{}

\begin{document}

\maketitle

% CORRECTED: Removed manual numbering
\section{The 2026 Vision and Its Central Execution Challenge}

% CORRECTED: Removed manual numbering
\subsection{Envision 2026: Hardening the Product, Planning for Change}

To support the strategy of market expansion and widespread adoption of the Envision system,
the Envision development team (TLRV-ML, KL) must resolve technical and classifier debt. This will result in a more robust and functional product. The Kumar Lab is proposing they work on the following three main deliverables:

\begin{enumerate}
    \item 24-hour classifier development cycle
    \item Continuous mHydraNet improvement through automatic keypoint labeling
    \item A first series of behavioral indices
\end{enumerate}

We see achievement of these three objectives critical to the 2026 Envision goals. Each of these deliverables will require changes to the Envision infrastructure and quick, continuous feedback amongst teams. The resolution of technical debt necessary to achieve the above deliverables will allow the Envision software infrastructure to produce higher quality data while also gaining resilience to product changes that facilitate adaptation to rapidly evolving market needs.

% CORRECTED: Removed manual numbering
\subsection{Mind the Gap: Areas of opportunity in the Envision system}

The success of current market is fundamentally and completely dependent on the machine learning (ML) algorithms that generate behavioral data. The 2026 Roadmap seeks to enter into new market— disposable cages, rat studies, and neuro-specific applications. These are all contingent on the underlying ML algorithms that generate the behavioral data, which is the core product Envision sells.

Currently, Envision is not reliably producing usable data. To produce valuable behavioral predictions that merits continued customer investment will require a monumental engineering effort that constitutes resolving a form of technical debt heretofore described as ”Classifier Debt”.

While technical debt refers to the accumulation of effort required to refactor the Envision infrastructure, Classifier Debt is the accumulated body of essential, non-negotiable redevelopment work required to make every existing digital measure in the Envision portfolio function accurately and reliably. This is not optional new feature development; it is a foundational prerequisite that must be paid down in full before any of the projected revenue or market share from these new segments can be realized. The scope of this debt is immense. Without the retrained algorithms, the new systems are unable to generate the data customers pay for at a baseline. The existence of this massive Classifier Debt creates a severe strategic risk. Every dollar of capital expenditure on developing and deploying new hardware solutions has a return on investment (ROI) of precisely zero until the Classifier Debt is fully paid down and the product can perform.

Paying down this debt is a large-scale undertaking that involves reviewing the current code infrastructure of the system and addressing its deficiencies. The KL has identified the following infrastructural deficiencies that must be addressed:

\textbf{1. MLOps Best Practices:}

(a) A sustainable solution for classifier development requires the adoption of industry-standard MLOps practices. This includes:

\begin{itemize}
    \item \textbf{Versioning and Traceability}: Every classifier, dataset, and training run must be version-controlled with full lineage tracking. This ensures that models can be reproduced, audited, and compared over time.
    \item \textbf{Automated Testing and Validation}: ML pipelines should include automated unit tests, integration tests, and validation checks at each stage (data ingestion, preprocessing, training, and deployment). Failures must halt the pipeline to prevent propagation of errors.
    \item \textbf{Continuous Integration/Continuous Deployment (CI/CD) for ML}: Classifiers should move seamlessly from development to staging to production using automated build-and-deploy pipelines, with guardrails for model performance thresholds.
    \item \textbf{Monitoring and Observability}: Production models must be continuously monitored for drift, degradation, and unexpected outputs. Automated alerts should flag when retraining or recalibration is required.
    \item \textbf{Data and Annotation Management}: Centralized tooling to manage annotations, training datasets, and metadata is essential to avoid duplication and maintain quality across teams.
    \item \textbf{Human-in-the-Loop (HITL) Integration}: The pipeline should explicitly support expert feedback loops, so annotations and error corrections feed directly back into the training cycle.
    \item \textbf{Scalability and Cloud-Native Infrastructure}: Pipelines should be containerized, orchestrated, and deployable across scalable cloud resources to ensure training and retraining can happen rapidly and reliably.
\end{itemize}
These best practices transform the Envision classifier pipeline into a repeatable, transparent, and continuously improving system, eliminating ad hoc workflows that currently create bottlenecks and strategic risk.

\textbf{2. Increased Test Coverage:}

Unit testing is currently insufficient and results in known instances of 'bit rot' code are still referred to in documentation. The code may or may not be designed with testing
in mind. Without test coverage, there is an increased risk of missing bugs and
subsequently releasing a non-functional product to users.

\textbf{3. Monolithic Development Approach}

(a) The current approach requires all iterations to be run through a large, complex,
and expensive infrastructure. Having to test small changes within the full system
is impossible without creating a manual work around, thus the design slows down
progress, punishes review, and prevents rapid prototyping and iteration.

\textbf{4. Opaque Outputs}

(a) The tooling to evaluate model outputs does not allow for effective assessment of
quality, qualitatively or quantitatively. This results in untraceable outputs, broken
feedback loops, decreased product quality, the inability to make evidence-based improvements,and finger-pointing amongst-teams.

% CORRECTED: Removed manual numbering
\subsection{The Long-Term Vision: A Platform Ecosystem}

This proposed integration should be viewed not as a bespoke solution for a single partner, but as the foundational pilot program for Envision’s long-term strategic ambition to become a true technology platform. The current roadmap outlines a future that includes an ”Integrated Development Environment,” public APIs, and a ”Digital Measure Marketplace” where customers and third parties can develop, host, and incorporate their own digital measures into an Envision study.

To build such a marketplace, Envision must first create the underlying infrastructure—the secure, scalable, and well-documented APIs—that allows an external development environment to seamlessly interact with the core platform. The work done to build the ”Optimized Data Flow” API for the Kumar Lab is precisely this foundational R\&D. Our partnership becomes the test case, the alpha user, and the first success story for this broader platform strategy. This reframes the engineering investment: it is not merely an operational expense to improve a single workflow, but a strategic R\&D investment in a future high-margin, platform-based revenue stream. It provides a dual return on investment by solving an immediate, critical bottleneck while simultaneously building the core infrastructure for a major future business initiative.

% CORRECTED: Removed manual numbering
\section{The Kumar Lab as a Strategic Partner}

The technical debt threatening the 2026 roadmap demands more than a simple augmentation of resources; it requires a fundamental shift in the operational paradigm of classifier development. KL proposes to collaborate more holistically with TLRV to mitigate risk to Envision’s success. As an aside, the KL recommends the Envision team go through the process of creating technical requirements using input from JMCRS’ product roadmap and strategy. The Kumar Lab will focus on three core deliverables that will drive improvements,reduce development time, and expand product functionality. They will each be outlined in detail below.

% CORRECTED: Removed manual numbering
\subsection{Detailed Proposals}

KL advocates for a philosophical change in product development. HydraNet and behavioral classifiers cannot be viewed as static, one-time deliverables. These are dynamic tools that must be continuously refined as the Envision platform encounters an ever-expanding distribution of data from new animals, environments, disease models, and hardware configurations. The roadmap, which implies the creation of dozens of distinct classifiers, necessitates a system built for evolution, not episodic development. Each of these ML Models must be treated as a living feature, subject to ongoing validation and improvement. 

The core challenge is not simply to build a classifier, but to build a sustainable process for its perpetual enhancement built on a robust ML Ops and testing framework built by TLRV. We propose a 24-hr classifier development cycle and a framework to rapidly retrain HydraNet that are outlined below. Finally, although behavior classifiers and a well functioning HydraNet are valuable commodity, the true value of Envision is in validated indices that allow users to make rapid decisions and tether them to the functionality of the Envision platform.


% CORRECTED: Removed manual numbering
\subsection{A 24-Hour Classifier Development Cycle}

To unlock the full potential of this strategic partnership and achieve the development velocity required by the 2026 roadmap, the underlying operational model for collaboration must be fundamentally re-architected. The current workflow is a significant source of friction and delay, making rapid, iterative development impossible. The proposed solution is to transition to an API-driven, integrated system designed to enable a ”24-hour classifier development cycle.”

% CORRECTED: Removed manual numbering
\subsubsection{How the API Unlocks the 24-Hour Behavior Classifier Development Cycle:}

The development of a high-performance behavioral classifier is an iterative scientific process, not a linear engineering task. A production-quality model is rarely achieved in a single pass; it typically requires three to five rounds of a rigorous ”human-in-the-loop” (HITL) cycle. This cycle involves training an initial model, running inference on new video data, having a human expert review the model’s predictions to identify failures and ”edge cases,” creating new annotations to correct these errors, and finally retraining the model with this enriched dataset.

This iterative refinement, which leverages human expertise to guide the algorithm, is essential for building robust and reliable models that can handle the complexities of real-world biological data.The strategic objective of the 24-hour development cycle is to compress each of these scientific iterations into a single day. This represents a paradigm shift to ”Agile Science,” where a researcher can test a hypothesis (a new model version) and receive actionable feedback for the next round of improvement within a 24-hour period. Achieving this velocity is impossible with the current ”hand-off” model but is the key to accelerating innovation and out-pacing competitors. To enable this agile workflow, a foundational technology bridge must be built in collaboration with TLRV, focused on three critical areas:

• \textbf{A Robust Publishing API:} This is the primary technical enabler. It automates the
submission of new model candidates from KL’s JABS environment into the Envision
platform, eliminating the manual integration bottleneck that currently introduces weeks of delay into each cycle. A well-designed API will support versioning, provide metadata rich responses (e.g., model version, performance metrics), and ensure scalable, secure integration between the two systems. To fully unlock the potential of an agile, distributed development cycle, the API must extend beyond model publishing to provide programmatic access to the underlying data ecosystem. Two critical capabilities are recommended:

\begin{itemize}
    \item \textbf{Study Metadata API}: The Kumar Lab must be able to query Envision metadata in a structured, automated fashion to identify the full diversity of animal and cage conditions available. This enables targeted sampling strategies for training datasets, ensuring classifiers are exposed to a broad distribution of environmental conditions and animal phenotypes. 
    
    \item \textbf{Artifact Access}: The API should also provide a secure mechanism for retrieving raw and processed artifacts (e.g., video segments and corresponding derived outputs such as pose and identity assignment). By specifying parameters such as cage ID, timestamp, pipeline ID, and artifact type, and other identifying attributes, API users can request short-lived pre-signed URLs that will facilitate integration into Kumar Lab’s on-premises annotation tools. This enables efficient workflows for expert annotation and model validation without duplicating or manually transferring large datasets nor does it require direct AWS access.

\end{itemize}Together, these additions extend the publishing API into a true data exchange interface, bridging Envision’s cloud infrastructure with the Kumar Lab’s local annotation and training environments. This ensures rapid, flexible, and secure access to the right data at the right time, further compressing iteration cycles and broadening the range of supported workflows.

•\textbf{ An Advanced MLOps Framework:} A mature MLOps pipeline which leverages industry-standard tooling is required to provide the backend automation for the 24-hour cycle. This includes capabilities for continuous and automated training workflows, comprehensive version control for code, data, and models, and scalable deployment to staging environments for rapid validation.

•\textbf{ Integrated Evaluation and Reporting Infrastructure:} To close the feedback loop, KL scientists must be able to efficiently review model performance directly within the
Envision application. This requires enhanced tooling for visualizing classifier outputs on video, comparing performance metrics between model versions, and flagging data segments for re-annotation. This infrastructure is the critical ”human-in-the-loop” interface that allows expert knowledge to be efficiently incorporated back into the model. In addition to publishing models into the Envision platform, the API should also enable the Kumar Lab to push new annotations and directly initiate model training in the cloud. This extension would allow the classifier development cycle to fully leverage distributed annotation workflows—ranging from a single expert annotator to larger, collaborative teams working asynchronously. By integrating annotation submission and cloud-based retraining triggers into the API, Envision can:

\begin{itemize}
    \item Eliminate bottlenecks caused by manual data transfer or coordination between environments.
    \item Provide flexible support for diverse research team structures, from individual subject matter experts to multi-annotator projects.
    \item Shorten the feedback loop by ensuring annotations immediately feed into the next training cycle.
    \item Further align the development process with the vision of an agile, iterative, and scalable pipeline for classifier evolution.

\end{itemize}

This approach not only accelerates the 24-hour development cycle but also ensures that the platform is future-proofed to support increasingly distributed and collaborative research efforts.

% CORRECTED: Removed manual numbering
\subsubsection{Deconstructing the Current Bottleneck: The ”Hand-Off” Model (Option 1)}

The current workflow where the ”KL team develops in JABS and ’hands-off’ model to Envision to be componentized,” represents the lowest and most inefficient level of integration. This ”throw it over the wall” approach is the direct root cause of the process issues the roadmap itself seeks to solve. The document notes that, regardless of the workflow, there is a pressing need for a clear ”Definition of ’done’,” ”Metrics to verify it is done,” and agreement on ”classifier definitions and labeling requirements”. These are not inherent challenges of ML development; they are symptoms of a disconnected process where one team works in isolation and hands a finished product to another, leading to costly rework when definitions, assumptions, or performance expectations are inevitably misaligned.

This model makes the Kumar Lab’s own stated goal of a ”24 hr classifier development loop” an operational impossibility. The administrative and engineering overhead of packaging a model, transferring it between environments, waiting for the TLRV team to write the componentization code, deploying it to a staging environment, and then running validation tests introduces weeks, if not months of latency into every single development cycle. Classifier development is an iterative scientific process, not a linear engineering one. The biological definition of a behavior like ”grooming” can be subtle and may require numerous refinements based on real-world data. The high latency of the current model makes this essential scientific iteration prohibitively slow and cumbersome.

% CORRECTED: Removed manual numbering
\subsubsection{The Pragmatic Solution: An API-Driven, ”Optimized Data Flow” (Option 2)}

The pragmatic and essential next step is the immediate implementation of the second option presented in the roadmap: to ”Optimize data flow between JABS and Envision such that models developed in JABS publish into Envision platform”. This model involves the creation of a robust Application Programming Interface (API) or a similar automated mechanism that allows the KL team to directly and independently push trained models, configuration files, and versioning metadata into the Envision ecosystem. The benefits of this approach are that it eliminates the manual ”hand-off” step, removing TLRV as a bottleneck in the development and validation cycle. This empowers the KL team to iterate and deploy new model versions independently and rapidly. This API-driven workflow creates the tight, high-fidelity feedback loop that is essential for refining complex behavioral classifiers. It allows KL scientists to see the output of their models running within the true Envision environment, using the platform’s own visualization tools like the Cage Dashboard to analyze performance. This directly supports Envision’s own goal to ”Enable smoother rerun of metrics on studies with data visualization in app” and is critical for ensuring models meet the stringent validation requirements of the scientific community.

% CORRECTED: Removed manual numbering
\subsection{A Novel Approach for Continual Evolution of mHydraNet}

The long-term scientific validity of every digital measure produced by the Envision platform rests upon the foundational ‘mHydraNet‘ model for animal detection, identification, and keypoint tracking. As Envision expands to new species, caging systems, and experimental contexts, this core model cannot be static; it requires a robust process for continuous improvement and validation. The current reliance on manual surface-based annotation for training data is a significant bottleneck that limits the speed, scale, and ultimate precision of model updates.

We propose using an ”Automatic key point labeling process” using ‘QD-Pi‘ (Quantum Dot based Pose estimation in vivo) in partnership with the Markowitz Lab. This method uses IR fluorescent dyes to label animal body parts. Animals are imaged under dual wavelength to rapidly create training data for mHydraNet. It provides a mechanism for rapid, semi-automated data annotation, which greatly reduces the time needed for a retraining loop while increasing label data by an order of magnitude. The strategic value of this methodology is twofold. First, it eliminates the ambiguity and inter-annotator variability inherent in manual labeling, which relies on subjective interpretation of surface features.

Second, it provides a mechanism to generate massive, perfectly labeled datasets for training and retraining ‘mHydraNet‘ with a level of accuracy that is unattainable with current methods. This enables a continuous improvement loop where, as new and challenging ”edge case” videos are collected, the ‘QD-Pi‘ system can be used to generate ground-truth data, ensuring ‘mHydraNet‘ becomes progressively more robust and accurate over time. The method is also species agnostic and will allow us to rapidly expand to rats and perhaps even other species. This collaborative work on foundational infrastructure also includes the parallel development of a more granular 12-keypoint pose estimation model.

Successfully operationalizing this cutting-edge technique requires a synergistic collaboration between KL and TLRV. While the Kumar Lab will provide the biomedical expertise for the nanoparticle injection protocols and data analysis, we require the engineering expertise of the TLRV team to construct and integrate the specialized data capture hardware into the Envision system. This joint effort will involve building multi-camera imaging arenas equipped with the specific NIR-optimized cameras, synchronized frame capture, synchronized excitation lighting (e.g., 660 nm LEDs), and
optical filters (e.g., 780 nm long-pass) necessary to reliably capture the nanoparticle fluorescence signal in a high-throughput manner. This initiative represents a foundational investment in the core quality and long-term viability of Envision’s entire data generation pipeline.

% CORRECTED: Removed manual numbering
\subsection{From Raw Data to Scientific Insight: The Power of Behavioral Indices}

The ultimate strategic value of this dynamic, agile approach is best exemplified by the ability to synthesize discrete digital measures into composite, scientifically validated indices. These indices become interpretable indicators of animal status, enabling users to make informed decisions.

The development of an automated, machine-vision-based Frailty Index (vFI) serves as a powerful example. Development of the vFI required over 800 animals from 8 weeks to 4.2 years old. Although burdensome to build, once completed and validated, behavioral  indices potentially lock the customer to a particular platform. A customer can move to different platforms to obtain a single behavior measure, but rebuilding and validating an entire index requires massive effort that customers are unlikely to do. In essence, it locks them into an ecosystem.

We should start to implement rudimentary indices as soon as possible. These can be as simple as combining grooming, scratching, and escape classifiers to build a ”stereotypy index,” or using feeding, drinking, and huddling to create a ”homeostatic behavior” index. This capability is directly aligned with Envision’s strategic need to offer validated, actionable insights, not just raw data streams. It provides a clear pathway to creating transformative, marketable endpoints—such as an ”ALS Progression Index” or a ”Neuropsychiatric Health Index”—that can serve as primary outcomes in preclinical research. This elevates the partnership from one focused on delivering features to one centered on co-creating the market-defining scientific products that will secure Envision’s leadership position.

% CORRECTED: Removed manual numbering
\subsubsection{Proposed Behavioral Indices}

% CORRECTED: Removed manual numbering
\paragraph{Stereotypy Index for Neuropsychiatric and Neurodevelopmental Models}

Stereotypy—the performance of excessive, repetitive, and invariant behaviors—is a core phenotype in models of autism, schizophrenia, and other neuropsychiatric disorders. We will develop a composite Stereotypy Index by first creating and validating a suite of classifiers for key component behaviors, including grooming, scratching, and escape behaviors. This index will be validated using the SSPsyGene (SING) mouse models, a large-scale NIMH consortium effort to characterize genes associated with psychiatric and neurodevelopmental disorders, in which stereotypy is a key behavioral domain. KL is funded to test over 100 neurodevelopmental disorder mouse lines using Envision.

% CORRECTED: Removed manual numbering
\paragraph{Neuromotor Index for Models of Motor Decline}

Progressive motor decline is a hallmark of numerous neurodegenerative disorders. To provide a sensitive and comprehensive measure of this process, we will develop a Neuromotor Index. This index will be built upon a foundation of classifiers for fundamental motor actions, including walking, rearing, turning, and climbing. The validation of this index will be conducted using a Rett Syndrome (RTT) mouse model, a severe neurodevelopmental disorder characterized by the regression of motor skills. KL has already established a collaborative framework for testing this model in JABS with Dr. Cat Lutz and the Rare Disease Translational Center (RDTC), providing a direct and efficient path for validation.

% CORRECTED: Removed manual numbering (and fixed the illogical 2.4.2.3)
\paragraph{Seizure 2.0 Index for Epilepsy and Convulsive Disorder Models}

To enhance Envision’s capabilities in epilepsy research, we will create a Seizure 2.0 Index designed to provide a more granular and quantitative assessment of seizure severity. This index will integrate classifiers for distinct components of the convulsive phenotype, including Straub tail, behavioral arrest, leg splaying, and tonic-clonic seizure events. The validation will be performed using a pentylenetetrazole (PTZ)-induced seizure model, a gold-standard pharmacological model for screening anticonvulsant compounds. By testing at both low and medium doses of PTZ, we can validate the index’s sensitivity across a range of seizure intensities.

% CORRECTED: Removed manual numbering
\subsubsection{A Note on Validation and Genetic Background}

A rigorous validation strategy is paramount to the scientific and commercial success of these indices. While the individual behavioral classifiers will be developed for robustness across multiple common genetic backgrounds, the final composite indices will be specifically validated in the context of the primary disease model’s background. For instance, the Neuromotor Index will be validated in the Rett model on its C57BL/6J background. For models like PTZ-induced seizures, which can be applied to multiple strains, validation across different genetic backgrounds will be pursued to demonstrate the broad applicability and robustness of the Seizure 2.0 Index.

% CORRECTED: Removed manual numbering
\begin{comment}
    

\subsection{Improved Performance, Robustness and Generalization of Upstream Models}

Upstream models—pose, segmentation, bounding box/offset prediction, and identity assignment—are foundational to downstream behavior classification. Improving their accuracy, robustness, and efficiency is critical to overall system performance.

\paragraph{Priorities}

\begin{itemize}
    \item Pose Estimation (12 Keypoints): Improve multi-mouse pose quality for DAX3 under cost/resource constraints. Enables finer posture modeling, reduces downstream noise, and supports integration with Kumar Lab gait analysis tools.
    \item Segmentation: Characterize current performance, then improve accuracy using a shared backbone with pose estimation for efficiency.
    \item Identity Assignment: Address frequent identity drops/swaps through new labels, architectural updates, and/or post-processing with pose, segmentation, and temporal context.
    \item Benchmarking: Compare mHydraNet against current state-of-the-art networks, evaluating both inference quality and computational efficiency to guide future model selection.
\end{itemize}

\paragraph{Cross-Cutting Considerations}

\begin{itemize}
    \item Model Selection: Favor models that balance predictive performance with resource efficiency, including scalable architectures or pretrained variants of varying complexity that can be tuned to specific use cases.
    \item Transferability: Prefer architectures adaptable to related tasks (segmentation, detection, identity).
    \item Fine-Tuning \& Community Support: Select models with accessible pretrained weights, well-documented adaptation pipelines, and active community adoption.
    \item Robustness: Validate across diverse recording conditions (lighting, enclosure geometry, camera angle) and scale variations to minimize retraining needs.
    \item Evaluation Framework: Define benchmarks for pose error, segmentation accuracy, identity continuity, robustness, and inference speed.
\end{itemize}

% CORRECTED: Removed manual numbering, changed to \subsubsection
\subsubsection{3D Pose Estimation}

Extending pose estimation from 2D to 3D will provide richer kinematic data and directly benefit downstream applications such as \textbf{gait analysis} and behavior classifiers for \textbf{rearing, grooming, and seizure detection}, which benefit depth cues or vertical posture information not captured by 2D. The strategy is to use synchronized multi-view cameras for \textbf{training data generation} (triangulated 3D ground truth), while deployed \textbf{inference will run on single top-down views}, learning to lift 2D keypoints into 3D space. This work is \textbf{complementary with Section 3.2}, where QD-Pi–based automatic labeling provides high-quality markers that can also be used to accelerate 3D ground-truth generation and continual downstream model improvement.

% CORRECTED: Removed stray \textbf{ and changed to \paragraph
\paragraph{Approach}

\begin{itemize}
    \item \textbf{Training Data:}
    \begin{itemize}
        \item Collect synchronized multi-view recordings to derive 3D ground truth via triangulation.
        \item Support both lights-on and IR conditions; apply augmentations to improve robustness.
        \item Optional depth sensors (e.g., ToF, IR projectors) may assist in generating labels but will not be part of deployed inference.
    \end{itemize}
    \item \textbf{Inference Strategy:}
    \begin{itemize}
        \item Train models to predict 3D pose from single-camera top-down inputs, supervised by multi-view 3D data.
        \item Explore temporal modeling (graph CNNs, transformer-based approaches) to leverage motion cues.
        \item 
    \end{itemize}
\end{itemize}

% CORRECTED: Removed manual numbering, changed to \paragraph, and removed extra }
\paragraph{Suggested Starting Points (SOTA Frameworks)}

\begin{itemize}
    \item \textbf{Proven 2D→3D Options:}
    \begin{itemize}
        \item \textbf{VideoPose3D} – established temporal convolutional lifting from 2D to 3D.
        \item \textbf{Anipose} – robust toolkit for multi-view calibration and 3D triangulation; ideal for generating training labels.
        \item \textbf{DeepLabCut 3D} – widely used in rodent labs; supports multi-view triangulation and 3D model training.
        \item \textbf{SLEAP (multi-view extension)} – supports multi-view triangulation; validated in 3D multi-animal scenarios.
    \end{itemize}
    \item \textbf{Strong 2D Backbones / Adaptable for 3D Pipelines:}
    \begin{itemize}
        \item \textbf{HRNet / HigherHRNet} – best-in-class 2D keypoint backbones; commonly paired with 3D lifting models.
        \item \textbf{STPoseNet / YOLO-Pose} – efficient one-stage 2D frameworks; not yet validated for 3D but promising as fast front-ends.
    \end{itemize}
\end{itemize}

% CORRECTED: Changed from bold text to \paragraph
\paragraph{Cross-Cutting Considerations}

\begin{itemize}
    \item \textbf{Resource Constraints:} Must run efficiently within Envision deployment requirements.
    \item \textbf{Robustness:} Models should generalize across cage geometries, lighting, occlusions, and scale variation.
\end{itemize}

% CORRECTED: Changed from bold text to \paragraph
\paragraph{Evaluation}
Define metrics including 3D joint position error, reprojection error, robustness under varied conditions, and compute/inference cost.

\end{comment}
\end{document}